%%%%%%%%%%%%%%%%%%%%%%%%%%%%%%%%%%%%%%%%%%%%%%%%%%%%%%%%%%%%%%%%%%%%%%
% Lớp 10 | Phương pháp tọa độ trong mặt phẳng (𝑂𝑥𝑦) | Đường thẳng trong mặt phẳng tọa độ
% Dạng 0H9?3-6: Bài toán tìm điểm
% Số câu: 18  (sinh tự động bởi tach_cau.py; từ nay thầy sửa trực tiếp file này)
%%%%%%%%%%%%%%%%%%%%%%%%%%%%%%%%%%%%%%%%%%%%%%%%%%%%%%%%%%%%%%%%%%%%%%

%%%%%%%%%%%%%%%% TRẮC NGHIỆM 4 PHƯƠNG ÁN %%%%%%%%%%%%%%%%

% nguồn: MaHoaNH_TracNghiem_Lop10
\begin{ex}%[0H9V3-6]%[ma:825408ab]
	Cho 3 điểm $A(-2;2)$, $B(7;5)$, $C(4;-5)$ và đường thẳng $\Delta\colon 2x+y-4=0$. Tìm tọa  độ điểm $N$ thuộc $\Delta$ sao cho $\left\vert\overrightarrow{NA}+\overrightarrow{NB}+\overrightarrow{NC}\right\vert $ đạt giá trị nhỏ nhất.
	 \choice
		 {$\left(\dfrac{29}{15};-\dfrac{2}{15}\right)$}
		 {\True $\left(\dfrac{29}{15};\dfrac{2}{15}\right)$}
		{$\left(-\dfrac{29}{15};\dfrac{2}{15}\right)$} 
		 {$\left(-\dfrac{7}{15};-\dfrac{1}{5}\right)$}
	 \loigiai{
		Gọi điểm $G$  thỏa $\overrightarrow{GA}+\overrightarrow{GB}+\overrightarrow{GC}=\overrightarrow{0}\Rightarrow \heva{&x_G=\dfrac{x_A+x_B+x_C}{3}=3\\&y_G=\dfrac{y_A+y_B+y_C}{3}=\dfrac{2}{3}}\Rightarrow G\left(3;\dfrac{2}{3}\right)$.\\
		Ta có $\left\vert\overrightarrow{NA}+\overrightarrow{NB}+\overrightarrow{NC}\right\vert = 3\vert\overrightarrow{NG}\vert=3NG$ nhỏ nhất $\Leftrightarrow GN\perp \Delta\Leftrightarrow N$  là hình chiếu của $G$ lên~ $\Delta$.\\
		Suy ra $N=d\cap  \Delta $ với $d$ qua $G$ và vuông góc với $\Delta$.\\
		$\Delta$ có một  vec-tơ pháp tuyến là $\overrightarrow{n}=(2;1)$.\\
		Suy ra $d$ có một vec-tơ chỉ phương là $\overrightarrow{u}=(2;1)\Rightarrow d$ có một  vec-tơ pháp tuyến là $\overrightarrow{n}_d=(1;-2)$.
		$$d\colon  1(x-3)-2\left(y-\dfrac{2}{3}\right)=0\Leftrightarrow x-2y-\dfrac{5}{3}=0.$$
		Tọa độ của  $N$ là nghiệm của hệ $\heva{&2x+y-4=0\\&x-2y-\dfrac{5}{3}=0}\heva{&x=\dfrac{29}{15}\\&y=\dfrac{2}{15}.}$\\
		Vậy $N\left(\dfrac{29}{15};\dfrac{2}{15}\right)$. 
	 }
\end{ex}


%%%%%%%%%%%%%%%% ĐÚNG/SAI %%%%%%%%%%%%%%%%

% nguồn: MaHoaNH_DungSai_Lop10
\begin{ex}%[0H9H3-6]%[ma:485b49c3]
	Trong mặt phẳng tọa độ $Oxy$, cho đường thẳng $d\colon \heva{&x=2-t\\& y=1+2t}$. Các  khẳng định sau đây, khẳng định nào đúng, khẳng định nào sai?
	 \choiceTF
		{\True Điểm $M(a;b)$  ($a<2$) thuộc đường thẳng $d$ sao cho $OM=\sqrt{10}$ thỏa mãn $3a-b=0$}
		{Gọi $N$ là điểm có hoành độ $c$ thuộc đường thẳng $d$ sao cho khoảng cách từ $N$ đến trục $Ox$ bằng  $3$. Tổng các giá trị của $c$ bằng $6$}
		{\True Điểm $P$ thuộc $d$ sao cho khoảng cách từ $A(1;2)$ đến $P$ nhỏ nhất có hoành độ là $\dfrac{7}{5}$} 
		{Điểm $Q$ thuộc $d$ sao cho $Q$  cách đều hai điểm $B(-1;2)$ và $C(2;-1)$ có tung độ là $3$}
	 \loigiai{
		\begin{itemchoice}
			\itemch Điểm $M\in  d$ nên $M(2-t;1+2t)\Rightarrow OM=\sqrt{(2-t)^2+(1+2t)^2}=\sqrt{5+5t^2}$.\\
			Suy ra $OM=\sqrt{10}\Leftrightarrow  5+5t^2=10\Leftrightarrow t=\pm 1$.\\
			Do $a<2$ nên $2-t<2 \Leftrightarrow t>0\Rightarrow t=1$, suy ra $M(1;3)$. Từ đó ta có $a=1$,  $b=3$, $3a-b=0$.
			\itemch Điểm $N\in d$ nên $N(2-t;1+2t)$. Khoảng cách từ $N$ đến $Ox$ bằng  $3$ nên
			$$|1+2t|=3 \Leftrightarrow \hoac{&t=1\\&  t=-2.}$$
			Suy ra $N(1;3)$  hoặc $N(4;-3)$, do đó $c\in \{1;4\}$.
			Vậy tổng  các giá trị của $c$ bằng $5$.
			\itemch Điểm  $P\in d$ nên $P(2-t;1+2t)$. Khi đó
			 $$AP=\sqrt{(2-t-1)^2+(1+2t-2)^2} = \sqrt{5t^2-6t+2} = \sqrt{5\left(t-\dfrac{3}{5}\right)^2+\dfrac{1}{5}} \geq \dfrac{\sqrt{5}}{5}.$$
			Vậy  $\min{AP}=\dfrac{\sqrt{5}}{5}$ khi $t=\dfrac{3}{5}$, khi đó $P\left(\dfrac{7}{5};\dfrac{11}{5}\right)$.
			\itemch Điểm $Q$  cách đều hai điểm $B$, $C$ nên $Q$ nằm trên đường trung trực của $BC$.\\
			Gọi $\Delta$ là đường trung trực của $BC$. Ta có $\overrightarrow{BC}=(3;-3)$  là véc-tơ pháp tuyến của $\Delta$, điểm $I\left(\dfrac{1}{2};\dfrac{1}{2}\right)$ là trung điểm của $BC$. Phương trình đường trung trực của $BC$ là $$\Delta \colon 3\left(x-\dfrac{1}{2}\right)-3\left(y-\dfrac{1}{2}\right)=0 \Leftrightarrow x-y=0.$$
			Tọa độ điểm $Q\in d\Rightarrow Q(2-t;1+2t)$,  mà $Q\in \Delta$ nên $2-t-1-2t=0 \Leftrightarrow t=\dfrac{1}{3}\Rightarrow Q\left(\dfrac{5}{3};\dfrac{5}{3}\right)$.
		\end{itemchoice}
	 }
\end{ex}

% nguồn: MaHoaNH_DungSai_Lop10
\begin{ex}%[0H9H3-6]%[ma:2a9c0c43]
	Trong mặt phẳng tọa độ $Oxy$, cho ba điểm $A(-2;2)$, $B(7;5)$, $C(4;-5)$ và đường thẳng $\Delta \colon 2x+y-4=0$. Các khẳng định sau đây, khẳng định nào đúng, khẳng định  nào sai?
	\choiceTF 
		{Điểm $M\in \Delta$ sao cho  $AM=3$ thì $M(1;2)$ hoặc $M\left(\dfrac{1}{5};\dfrac{18}{5}\right)$}
		{\True  Điểm $N(x_N;y_N)\in \Delta$ sao cho $BN$ ngắn nhất thỏa mãn $x_N-y_N=-1$}
		{\True Điểm $P\in  \Delta$ và cách đều hai điểm $A$, $B$ có tung độ bằng $-10$}
		{Điểm $Q\in \Delta$ sao  cho $\left|\overrightarrow{QB}+\overrightarrow{QC}\right|$ ngắn nhất có hoành độ bằng $-\dfrac{27}{5}$}
	 \loigiai{
		\begin{itemchoice}
			\itemch Vì điểm  $M\in \Delta$ nên $M(m;4-2m)$, suy ra
			$$AM=3 \Leftrightarrow (m+2)^2+(2-2m)^2=9 \Leftrightarrow 5m^2-4m-1=0 \Leftrightarrow \hoac{&m=1\\& m=-\dfrac{1}{5}.}$$ 
			Vậy $M(1;2)$  hoặc $ M\left(-\dfrac{1}{5};\dfrac{22}{5}\right)$.
			\itemch Điểm $N\in \Delta$ sao cho $BN$ ngắn nhất thì $N$ là hình chiếu vuông  góc của $B$ trên $\Delta$. Khi đó $N$ thuộc đường thẳng $d$ đi qua $B$ và vuông góc với $\Delta$.\\
			Phương trình đường thẳng  $d\colon (x-7)-2(y-5)=0 \Leftrightarrow x-2y+3=0$.\\
			Tọa độ  điểm $N$ là nghiệm của hệ $\heva{&2x+y-4=0\\& x-2y+3=0} \Leftrightarrow \heva{&x=1\\& y=2} \Rightarrow N(1;2)$.\\
			Suy ra $x_N=1$, $y_N=2$ và  $x_N-y_N=-1$.
			\itemch Điểm $P$ cách đều hai điểm $A$, $B$ nên nằm trên đường trung trực của $AB$. Gọi $d_1$ là đường trung trực của $AB$, ta có $\overrightarrow{AB}=(9;3)$ là véc-tơ pháp tuyến của $d_1$,  điểm $I\left(\dfrac{5}{2};\dfrac{7}{2}\right)$ là trung điểm của $AB$, suy ra phương trình $d_1$ là
			 $$9\left(x-\dfrac{5}{2}\right)+3\left(y-\dfrac{7}{2}\right)=0 \Leftrightarrow 3x+y-11=0.$$
			Tọa độ $P$ là nghiệm của hệ $\heva{&2x+y-4=0\\& 3x+y-11=0} \Leftrightarrow \heva{&x=7\\& y=-10}\Rightarrow  P(7;-10)$.
			\itemch Vì $Q\in \Delta$  nên $Q(a;4-2a)$. Gọi $D$ là trung điểm của $BC$, ta có $D\left(\dfrac{11}{2};0\right)$ và $\overrightarrow{QB}+\overrightarrow{QC} = 2\overrightarrow{QD}$, suy ra
			 \allowdisplaybreaks
			\begin{eqnarray*}
				\left|\overrightarrow{QB} +\overrightarrow{QC}\right| = 2QD &=& 2\sqrt{\left(\dfrac{11}{2}-a\right)^2 +(2a-4)^2} =  2\sqrt{5a^2-27a+\dfrac{185}{4}}\\
				&=& 2\sqrt{5\left(a-\dfrac{27}{10}\right)^2+\dfrac{49}{5}} \geq \dfrac{14\sqrt{5}}{5}. 
			\end{eqnarray*}
			Vậy $\min\left|\overrightarrow{QB} +\overrightarrow{QC}\right| = \dfrac{14\sqrt{5}}{5}$  khi $a=\dfrac{27}{10}$, khi đó điểm $Q$ có hoành độ bằng $\dfrac{27}{10}$.
		\end{itemchoice}
	} 
\end{ex}

% nguồn: MaHoaNH_DungSai_Lop10
\begin{ex}%[0H9V3-6]%[ma:61da3251]
	Trong mặt phẳng $(Oxy)$, cho điểm $M(2;4)$ và  đường thẳng $d\colon \heva{&x=1-3t\\&y=2+t}$ ($t\in\mathbb{R}$).
	\choiceTF 
		{\True Tọa độ véc-tơ  $\overrightarrow{OM}$ là $(2;4)$}
		{Véc-tơ $\overrightarrow{n}=(-3;1)$ là một véc-tơ pháp tuyến của  đường thẳng $d$}
		{\True Đường thẳng đi qua $M$ và vuông góc với đường thẳng  $d$ có phương trình là $3x-y-2=0$}
		{Gọi  $A(a;b)$ là điểm nằm trên đường thẳng $d$ sao cho $MA$ nhỏ nhất. Khi đó $a+b=3$}
	 \loigiai{
		Đường thẳng $d$ có véc-tơ chỉ phương $\overrightarrow{u}=(-3;1)$. 
		\begin{itemchoice}
			\itemch Vì $M(2;4)$  nên $\overrightarrow{OM}=(2;4)$.
			\itemch Véc-tơ $\overrightarrow{n}=(-3;1)$ là một véc-tơ chỉ phương  của đường thẳng $d$.
			\itemch Đường thẳng $\Delta$ qua $M$ và vuông góc với đường thẳng $d$ nhận véc-tơ $\overrightarrow{u}=(-3;1)$ làm véc-tơ pháp tuyến nên có phương trình  là
			$$-3(x-2)+(y-4)=0 \Leftrightarrow  3x-y-2=0.$$
			\itemch Điểm $A$ nằm trên đường thẳng $d$ nên $A(1-3t;2+t)$. Khi đó $MA$ nhỏ nhất khi $A$ là hình chiếu vuông góc của $M$ trên $d$, suy ra $A\in \Delta \Rightarrow 3(1-3t)-(2+t)-2=0  \Leftrightarrow t=-\dfrac{1}{10} \Rightarrow A\left(\dfrac{13}{10}; \dfrac{19}{10}\right)$.\\
			Vậy $a=\dfrac{13}{10}$, $b=\dfrac{19}{10}$ và $a+b=\dfrac{32}{10}$. 
	\end{itemchoice}}
\end{ex}

% nguồn: MaHoaNH_DungSai_Lop10
\begin{ex}%[0H9V3-6]%[ma:480d3e94]
	Trong mặt phẳng tọa độ $Oxy$, cho hai đường thẳng $\Delta_1\colon  \heva{&x=1+3t \\
	&y=3-4t 
	}$ và  $\break \Delta_2\colon 4x-3y-5=0$.
	 \choiceTF
		{Hai đường thẳng $\Delta_1$ và $\Delta_2$ vuông  góc với nhau}
		{\True Khoảng cách từ điểm  $O(0; 0)$ đến đường thẳng $\Delta_2$ bằng $1$}
		{Điểm  $A(-1;-3)$ thuộc đường thẳng $\Delta_1$}
		{Điểm $M\left(\dfrac{9}{4}; \dfrac{3}{4}\right)$ là tọa  độ giao điểm của hai đường thẳng $\Delta_1$ và $\Delta_2$}
	 \loigiai{
		\begin{itemchoice}
			\itemch Xét đường thẳng $\Delta_1$  có một vectơ pháp tuyến $\overrightarrow{n}_1=(4; 3)$.\\
			Xét đường thẳng $\Delta_2$  có một véc-tơ pháp tuyến $\overrightarrow{n}_2=(4;-3)$.\\
			Vì $\overrightarrow{n_1} \cdot \overrightarrow{n_2}=4.4+3(-3)=7\neq 0$ nên hai đường thẳng $\Delta_1$ và $\Delta_2$ không vuông  góc với nhau.
			\itemch Ta có $\mathrm{d}\left(O;  \Delta_2\right)=\dfrac{|4\cdot0-3\cdot0-5|}{\sqrt{4^2+(-3)^2}}=1$.
			 \itemch Xét $\heva{&-1=1+3t \\
			&-3=3-4t} \Leftrightarrow\heva{&t=-\dfrac{2}{3} \\ 
			 &t=\dfrac{2}{3}
			}$(vô  lý).\\
			Vậy  $A\notin \Delta_1$.
			\itemch Đường thẳng $\Delta_1$ đi qua điểm $B(1; 3)$ và có một vectơ pháp tuyến $\overrightarrow{n}_1=(4; 3)$  có phương trình $4(x-1)+3(y-3)=0\Leftrightarrow 4x+3y-13=0$.\\
			Xét hệ phương  trình $\heva{&4x+3y-13=0\\
			 &4x-3y-5=0} \Leftrightarrow\heva{&x=\dfrac{9}{4} \\
			&y=\dfrac{4}{3}.}$\\ 
			Vậy $M\left(\dfrac{9}{4}; \dfrac{4}{3}\right)$ là tọa độ giao điểm của hai đường thẳng $\Delta_1$ và  $\Delta_2$.
	\end{itemchoice}}
\end{ex}

% nguồn: MaHoaNH_DungSai_Lop10
\begin{ex}%[0H9V3-6]%[ma:5ef43ed0]
	Trong mặt phẳng tọa độ $Oxy$, cho tam giác $ABC$ có $A(1;-2)$ và đường thẳng chứa cạnh $BC$ có phương trình $5x-3y+1=0$. $K$ là một điểm nằm trên đoạn  thẳng $AH$ sao cho $\overrightarrow{AK}=\dfrac{3}{4} \overrightarrow{AH}$
	\choiceTF 
		{\True Một vectơ chỉ phương của đường thẳng $BC$ là $\vec{u}_{BC}=(3;  5)$}
		{\True  Đường cao $AH$ có phương trình là $3x+5y+7=0$}
		{Hoành độ của  điểm $H$ là một số nguyên dương}
		{Có hai điểm $K$ thỏa  mãn yêu cầu bài toán}
	\loigiai{ 
		\begin{itemchoice}
			\itemch 
			Đường thẳng chứa cạnh $BC$ có phương trình $5x-3y+1=0$, suy  ra nó có một vectơ chỉ phương $\vec{u}_{BC}=(3; 5)$.
			\itemch 
			Đường cao  $AH$ đi qua điểm $A(1;-2)$ và vuông góc với đường thẳng $BC$ nên có vectơ pháp tuyến là $\vec{n}_{AH}=\vec{u}_{BC}=(3; 5)$.\\
			Do đó phương trình đường cao  $AH$ là $3(x-1)+5(y+2)=0\Leftrightarrow 3x+5y+7=0$.
			\itemch Vì $\{H\}=AH\cap BC$ suy ra tọa độ của $H$ là nghiệm của  hệ phương trình
			$$\heva{&3x+5y+7=0\\ 
			&5x-3y+1=0} \Leftrightarrow\heva{&3x+5y=-7\\ 
			&5x-3y=-1} \Leftrightarrow\heva{&x=-\dfrac{13}{17} \\ 
			 &y=-\dfrac{16}{17}}$$
			 suy ra $H\left(-\dfrac{13}{17};-\dfrac{16}{17}\right)$.
			\itemch 
			Giả sử $K(x;  y)$ nên $\overrightarrow{AK}=(x-1; y+2)$, $\overrightarrow{AH}=\left(-\dfrac{13}{17}-1;-\dfrac{16}{17}+2\right)$.\\
			Nên $\dfrac{3}{4} \overrightarrow{AH}=\left(-\dfrac{90}{68}; \dfrac{54}{68}\right) \Rightarrow \dfrac{3}{4} \overrightarrow{AH}=\left(-\dfrac{45}{34};  \dfrac{27}{34}\right)$.\\
			Giả thiết $\overrightarrow{AK}=\dfrac{3}{4} \overrightarrow{AH}$  suy ra
			$\heva{&x-1=-\dfrac{45}{34} \\ 
			&y+2=\dfrac{27}{34}} \Leftrightarrow\heva{&x=-\dfrac{11}{34} \\ 
			 &y=-\dfrac{41}{34}.}$\\
			Vậy  $K\left(-\dfrac{11}{34};-\dfrac{41}{34}\right)$.
	\end{itemchoice}}
\end{ex}

% nguồn: MaHoaNH_DungSai_Lop10
\begin{ex}%[0H9V3-6]%[ma:22fc7e71]
	Trong mặt phẳng  $Oxy$, cho hai điểm $A(1;2)$, $B(3;-1)$.
	 \choiceTF
		{\True Một véctơ pháp  tuyến của đường thẳng $AB$ là $\vec{n}=(3;2)$}
		{\True Điểm $M$ trên trục $Ox$ sao cho $MA+MB$  nhỏ nhất có tọa độ $\left(\dfrac{7}{3}; 0\right)$}
		{Phương trình tổng quát của đường trung trực của đoạn thẳng $AB$  là $2x-3y-1=0$}
		{Đường thẳng $AB$ vuông  với đường thẳng $d\colon 3x-2y+2024=0$}
	 \loigiai{
	\begin{itemize}
		\item Ta có  $\vec{AB}=(2;-3)$ là một véc-tơ chỉ phương của đường thẳng $AB$.\\
		 Do đó đường thắng $AB$ có một véc-tơ pháp tuyến là $\vec{n}=(3;2)$.
		\item Đường thẳng trung trực của đoạn thẳng $AB$ đi qua trung điểm $I\left( 2; \dfrac{1}{2}\right)$,  nhận $\vec{AB}=(2;-3)$ làm véc-tơ pháp tuyến, có phương trình là
		\[ 
		2(x-2) - 3\left( y-\dfrac{1}{2}\right)=0 \Leftrightarrow  2x-3y-\dfrac{5}{2}=0.
		 \]
		\item Đường  thẳng $d$ có véc-tơ pháp tuyến là $\vec{n}_1= (3;-2)$.\\
		Ta thấy $\dfrac{2}{3}\neq \dfrac{-3}{-2}$ nên hai véc-tơ $\vec{AB}$ và $\vec{n}_1$ không cùng phương. Do đó đường thẳng $AB$ không vuông góc góc với đường  thẳng $d$.
		\item Ta thấy hai điểm $A(1;2)$, $B(3;-1)$ ở khác phía  so với trục $Ox$, do $y_A=2>0$ và $y_B=-1<0$.\\
		Gọi $M(m;0)$ là điểm trên trục $Ox$.  Khi đó $MA+MB\ge AB$. Dấu ``='' xảy ra khi $M$ là giao điểm của $AB$ và trục $Ox$, tức là ba điểm $A$, $M$, $B$ thẳng hàng.\\
		 Ta có $\vec{AB}=(2;-3)$, $\vec{AM}=(m-1;-2)$.
		\begin{eqnarray*}
			\text{Ba điểm } A, M, B \text{ thẳng hàng} &\Leftrightarrow& \text{hai véc-tơ } \vec{AB}, \vec{AM} \text{  cùng phương}\\
			 &\Leftrightarrow& \dfrac{m-1}{2}= \dfrac{-2}{-3} \Leftrightarrow m=\dfrac{7}{3}.
		\end{eqnarray*}
		Vậy toạ  độ $M\left(\dfrac{7}{3}; 0\right)$.
	\end{itemize}
	\begin{itemchoice}
		\itemch Đúng. 
		\itemch Đúng. 
		\itemch  Sai.
		 \itemch Sai.
	\end{itemchoice}
	} 
\end{ex}


%%%%%%%%%%%%%%%% TRẢ LỜI NGẮN %%%%%%%%%%%%%%%%

% nguồn: MaHoaNH_TraLoiNgan_Lop10
\begin{bt}%[0H9H3-6]%[ma:3e113595]
	Trong mặt phẳng $Oxy$, cho đường thẳng $d\colon x-2y+1=0$ và điểm $M(2;-2)$. Tọa độ hình chiếu vuông góc  của điểm $M$ lên đường thẳng $d$ là $N(a;b)$. Khi đó $a\cdot b$ bằng bao nhiêu? (viết kết quả dưới dạng số thập phân).
	\shortans{0,48} 
	\loigiai{ 
		Đường thẳng $d$ có một  vectơ pháp tuyến là $\vec{n}_d=(1;-2)$.\\
		Suy ra vectơ  chỉ phương của $d$ là $\vec{u}_d=(2;1)$.\\
		Gọi $d'$ là đường thẳng đi qua $M$ và vuông góc với $d$, khi đó $d'$ nhận  vectơ chỉ phương của $d$ làm một vectơ pháp tuyến.\\
		Suy ra vecto pháp  tuyến của $d'$ là $\vec{n}_{d'}=(2;1)$.\\
		Phương trình đường thẳng $d'$ là $2(x-2)+(y+2)=0\Leftrightarrow  2x+y-2=0$.\\
		Gọi $N$ là giao điểm của $d$ và $d'$,  tọa độ điểm $N$ là nghiệm của hệ phương trình
		$$\heva{&x-2y=-1\\&2x+y=2}\Leftrightarrow  \heva{&x=\dfrac{3}{5}\\&y=\dfrac{4}{5}.}$$
		 Vậy hình chiếu vuông góc của $M$ lên đường thẳng $d$ là $N\left(\dfrac{3}{5};\dfrac{4}{5}\right)$.\\
		Suy ra $a\cdot  b=\dfrac{12}{25}=0{,}48$.
	 }
\end{bt}

% nguồn: MaHoaNH_TraLoiNgan_Lop10
\begin{bt}%[0H9H3-6]%[ma:bf211799]
	Cho $\triangle ABC$ có trọng tâm $G(-2;-1)$; $AB\colon 4x+y+15=0$; $AC\colon 2x+5y+3=0$. Gọi $C(x_0;y_0)$, khi  đó $x_0+2y_0$ bằng bao nhiêu?
	 \shortans{-1}
	 \loigiai{
		Tọa độ điểm $A=AB\cap AC$ là  nghiệm của
		 hệ $\heva{&{4x+y+15=0} \\&{2x+5y+3=0}}\Leftrightarrow \heva{&{x=-4} \\ &{y=1}}\Rightarrow A(-4;1).$
		Ta  có
		\begin{itemize}
			\item $B\in  AB\colon y=-4x-15\Rightarrow B\left(x_B;-4x_B-15\right)$.
			 \item $C\in AC:y=\dfrac{-2x-3}{5} \Rightarrow C\left(x_C;\dfrac{-2x_C-3}{5} \right)$.
		\end{itemize}
		Vì $G$ là trọng tâm  $\triangle ABC\Leftrightarrow \heva{&{x_A+x_B+x_C=3x_G} \\
		&{y_A+y_B+y_C=3y_G}}\Leftrightarrow  \heva{&{-4+x_B+x_C=3(-2)} \\
		&{-1-4x_B-15+\dfrac{-2x_C-3}{5}=-3}}\\\Leftrightarrow \heva{&{x_B=-3}  \\&{x_C=1}}\Rightarrow B(-3;-3),C(1;-1)$.\\
		Suy ra,  $x_0=1$ và $y_0=-1$ do đó, $x_0+2y_0=1-2=-1$.
	} 
\end{bt}

% nguồn: MaHoaNH_TraLoiNgan_Lop10
\begin{bt}%[0H9H3-6]%[ma:3c0ce9cd]
	Trong mặt phẳng toạ độ $Oxy$, cho hình thoi $ABCD$ có $A(0;2),B(4;3)$, giao điểm hai đường chéo nằm trên đường  thẳng $\Delta:x-3y=0$. Biết tọa độ đỉnh $D(x_0;y_0)$ có $x_0<0$. Tính $x_0+y_0$.
	 \shortans{-3}
	\loigiai{ 
		\begin{center}
			\begin{tikzpicture}
					 \def\canh{4}
					\coordinate (A) at (0,0); 
					\coordinate (B) at ($(A)+(-65:\canh)$); 
					 \coordinate (D) at ($(A)+(-115:\canh)$);
					\coordinate (C) at  ($(B)+(D)-(A)$);
					\draw (A)--(B)--(C)--(D)--cycle (A)--(C) 
						(B)--(D); 
					\coordinate (I) at (intersection of A--C and B--D); 
					\foreach  \i/\g in {A/90,B/0,C/-90,D/180,I/45}{\draw[fill=white](\i) circle (1pt)
						($(\i)+(\g:3mm)$)  node[scale=1]{$\i$};}
				\end{tikzpicture}
		\end{center}
		Gọi $I$ là giao điểm của hai đường chéo. Vì $I$ thuộc $\Delta$  nên giả sử $I(3t;t)$.\\
		Khi đó  $\overrightarrow{IA}=(-3t;2-t),\overrightarrow{IB}=(4-3t;3-t)$.\\
		 Vì tứ giác $ABCD$ là hình thoi nên
		\[\overrightarrow{IA}\cdot  \overrightarrow{IB}=0\Leftrightarrow (-3t)(4-3t)+(2-t)(3-t)=0\Leftrightarrow 10t^2-17t+6=0.\]
		Suy  ra $t=\dfrac{1}{2}$ hoặc $t=\dfrac{6}{5}$.
		\begin{itemize}
			 \item Với $t=\dfrac{1}{2}$ ta có: $I\left(\dfrac{3}{2};\dfrac{1}{2} \right)\Rightarrow C(3;-1),D(-1;-2)$ (nhận).\\
			Suy  ra, $x_0+y_0=-1-2=-3$.
			\item Với $t=\dfrac{6}{5}$ ta có: $I\left(\dfrac{18}{5};\dfrac{6}{5} \right)\Rightarrow C\left(\dfrac{36}{5};\dfrac{2}{5}  \right),D\left(\dfrac{16}{5};-\dfrac{3}{5} \right)$ (loại).
		\end{itemize}
	} 
\end{bt}

% nguồn: MaHoaNH_TraLoiNgan_Lop10
\begin{bt}%[0H9V3-6]%[ma:6c602f7f]
	Trong mặt phẳng $Oxy$, cho đường thẳng $d\colon x-2y+1=0$ và điểm $M(2;-2)$. Tọa độ hình chiếu vuông góc của điểm $M$  lên đường thẳng $d$ là $N(a; b)$. Khi đó $ab$ bằng bao nhiêu?
	\shortans{0,48} 
	\loigiai{ 
		Đường thẳng $d$ có một vectơ pháp tuyến là $\overrightarrow{n_d}=(1;-2)$. Suy ra vectơ chỉ phương của $d$ là $\overrightarrow{u_d}=(2; 1)$.\\ 
		Gọi $d'$ là đường thẳng đi qua $M$ và vuông góc  với $d$, khi đó $d'$ nhận vectơ chỉ phương của $d$ làm một vectơ pháp tuyến, do đó, $\overrightarrow{n_{d'}}=(2; 1)$.\\
		Phương trình đường  thẳng $d'$ là $2(x-2)+(y+2)=0\Leftrightarrow 2x+y-2=0$.\\
		Gọi $N$ là giao điểm của $d$ và $d'$, tọa độ điểm $N$ là nghiệm của  hệ $\heva{&x-2y=-1 \\&2x+y=2}\Leftrightarrow\heva{&x=\dfrac{3}{5} \\&y=\dfrac{4}{5}.}$\\
		Vậy hình chiếu vuông góc  của $M$ lên đường thẳng $d$ là $N\left(\dfrac{3}{5}; \dfrac{4}{5}\right) \Rightarrow a \cdot b=\dfrac{12}{25}=0{,}48$.
	 }
\end{bt}

% nguồn: MaHoaNH_TraLoiNgan_Lop10
\begin{bt}%[0H9V3-6]%[ma:c15db22a]
	Cho $\triangle ABC$ có trung điểm cạnh $BC$  là $M(-1,-1)$; $AB\colon x+y-2=0$; $AC\colon 2x+6y+3=0$. Tính diện tích $ABC$.
	 \shortans{10}
	\loigiai{ 
		Tọa độ điểm $A=AB\cap AC$ là nghiệm của hệ $\left\{\begin{array}{l}{x+y-2=0} \\ {2x+6y+3=0}
	\end{array}\Leftrightarrow \left\{\begin{array}{l}{x=\dfrac{15}{4}} \\ {y=\dfrac{-7}{4}}
	\end{array}\Rightarrow A\left(\dfrac{15}{4};\dfrac{-7}{4} \right).\right. \right.$
	 Ta có
	\begin{itemize}
		\item $B\in AB\colon y=-x+2\Rightarrow B\left(x_B;-x_B+2\right)$. 
		\item $C\in AC\colon y=\dfrac{-2x-3}{6} \Rightarrow  C\left(x_c;\dfrac{-2x_c-3}{6} \right)$.
	\end{itemize}
	Vì $M$ là trung điểm của $BC\Leftrightarrow \left\{\begin{array}{l}{x_B+x_C=2x_M} \\ {y_B+y_C=2y_M}
	\end{array}\Leftrightarrow \left\{\begin{array}{l}{x_B+x_C=-2} \\ {-x_B+2+\dfrac{-2x_C-3}{6}=-2}
	\end{array}\right. \right.$
	\[\Leftrightarrow \left\{\begin{array}{l}{\begin{array}{l}{x_B+x_C=-26} \\ {x_B-2x_C=-21}
	\end{array}}
	\end{array}\Leftrightarrow \left\{\begin{array}{l}{\begin{array}{l}{x_B=\dfrac{25}{4}} \\ {x_C=\dfrac{-33}{4}}
	\end{array}}
	\end{array}\Rightarrow B\left(\dfrac{25}{4};\dfrac{-17}{4} \right),C\left(\dfrac{-33}{4};\dfrac{9}{4} \right).\right. \right.
	 \]
	Ta có $\overrightarrow{AB}=\left(\dfrac{5}{2};-\dfrac{5}{2}\right)$ và  $\overrightarrow{AC}=\left(-12;4\right)$.\\
	Khi  đó, diện tích $ABC$ là $S_{ABC}=\dfrac{1}{2}\cdot \left|\dfrac{5}{2}\cdot 4-\left(-\dfrac{5}{2}\right)\cdot (-12)\right|=10$.
	} 
\end{bt}

% nguồn: MaHoaNH_TraLoiNgan_Lop10
\begin{bt}%[0H9C3-6]%[ma:c62d0af4]
	Trong mặt phẳng với hệ tọa độ $Oxy$, cho tam giác $ABC$ vuông ở $A$. Biết $A(-1;4)$, $B(4;-1)$, đường thẳng  $BC$ đi qua điểm $K\left(\dfrac{7}{3};2\right)$. Biết tọa độ đỉnh $C(x_C;y_C)$, tính $y_C-x_C$.
	\shortans{5} 
	 \loigiai{
		Đường thẳng $BC$ có một  véc-tơ chỉ phương là $\overrightarrow{BK}=\left(-\dfrac{5}{3};3\right)$ và đi qua điểm $B(4;-1)$ nên có phương trình là
		$$BC\colon 3(x-4)+\dfrac{5}{3}(y+1)=0\Leftrightarrow 9x+5y-31=0.$$ 
		Vì $C\in BC$ nên $C\left(a;\dfrac{31-9a}{5}\right)$.  Khi đó
		$$\overrightarrow{AB}=(5;-5),\overrightarrow{AC}=\left(a+1;\dfrac{11-9a}{5}\right).$$ 
		Tam giác $ABC$ vuông ở $A$  nên $AB\perp AC$ do đó
		$$\overrightarrow{AB}\cdot\overrightarrow{AC}=0\Leftrightarrow  5(a+1)-5\cdot\dfrac{11-9a}{5}=0\Leftrightarrow 14a-6=0\Leftrightarrow a=\dfrac{3}{7}\Rightarrow C\left(\dfrac{3}{7};\dfrac{38}{7}\right).$$
		Suy ra  $y_C-x_C=\dfrac{38}{7}-\dfrac{3}{7}=5$.
	 }
\end{bt}

% nguồn: MaHoaNH_TraLoiNgan_Lop10
\begin{bt}%[0H9C3-6]%[ma:22054c07]
	Trên hệ trục tọa độ $Oxy$, cho hai điểm $A(3;1)$, $B(0;1)$ và các đường thẳng $\Delta_1\colon 2x-y-2=0$, $\Delta_2\colon \heva{&x=-3+t\\&y=1-t}$ ($t$ là tham số). Gọi $M$ và $N$ lần lượt là các điểm thay đổi trên $\Delta_1$ và $\Delta_2$. Biết giá trị nhỏ nhất $T_{\min}$ của biểu thức $T=AM+BN+MN$ có dạng $T_{\min} =  a\sqrt{b}$ với $a$ là số nguyên dương, $b$ là số nguyên tố. Tính $a^2+2b$.
	 \shortans{19}
	 \loigiai{
	\immini{ 
		Gọi $B'$ là điểm đối xứng của  $B$ qua $\Delta_2$.\\
		Gọi $M'$, $N'$ lần lượt là giao điểm của $AB'$ với $\Delta_1$,  $\Delta_2$.\\
		Ta có  $NB=NB'$.\\
		$T=AM+BN+MN=AM+B'N+MN\geq  AB'$.\\
		Suy ra $T_{\min}=AB'$ khi đó $M\equiv M'$,  $N\equiv N'$.\\
		Gọi $d$ là đường thẳng qua  $B$ và vuông góc $\Delta_2$\\
		 $\Delta_2\colon \heva{&x=-3+t\\&y=1-t}\Leftrightarrow x+y+2=0$.\\
		Suy ra  $d\colon x-y+m=0$.
		 }{
		\begin{tikzpicture}[line join=round,line cap=round,font=\footnotesize,scale=1,>=stealth]
				 \path
					(3,1)  coordinate (A)
					(0,1)  coordinate (B)
					 (-2,0) coordinate (C)
					(0,-2) coordinate  (D)
					(-3,1) coordinate  (N)
					(2,2) coordinate (M) 
					(1,0)  coordinate (M')
					($(C)!(B)!(D)$) coordinate  (H)
					($(B)!2!(H)$)  coordinate (B')
					(intersection  of C--D and A--B') coordinate (N');
				\draw[-stealth](-3.3,0)--(3.3,0)  node[above]{$x$};
				\draw[-stealth](0,-3)--(0,3)  node[right]{$y$};
				\draw[smooth,samples=300,domain=-0.5:2.4]  plot (\x,{2*(\x)-2}) node[right]{$\Delta_1$};
				\draw[smooth,samples=300,domain=-3.2:0.5] plot (\x,{-(\x)-2}) node[right]{$\Delta_2$}; 
				\draw  (B)--(B')--(A)
					 (B)--(N)--(B')
					 (A)--(M)--(N);
				\fill (0,0) circle (1pt) node[below  left]{$O$};
				 \draw[dashed](3,0)|-(0,1);
				 \draw pic[draw,angle radius=5pt]{right angle=B--H--N};
				\foreach \x/\g in {-2/-120,1/-80,3/-90}\fill[black] (\x,0)  circle (1pt)+(\g:.3) node{$\x$};
				\foreach \x/\g in {-2/180,1/60}\fill[black] (0,\x)  circle (1pt)+(\g:.4) node{$\x$};
				\foreach \x/\g in  {A/60,B/140,B'/180,N/90,N'/-90,M/140,M'/120,H/0}\fill[black] (\x) circle (1pt)+(\g:.4) node{$\x$};
			\end{tikzpicture}
		} 
		 \noindent $B(0;1)\in d\Rightarrow 0-1+m=0\Leftrightarrow m=1$.\\
		$d\colon x-y-1=0$.\\ 
		Vì $H=d\cap\Delta_2$ nên tọa độ của $H$  là nghiệm của hệ
		 $$\heva{&x+y+2=0\\&x-y-1=0}\Leftrightarrow \heva{&x=-\dfrac{3}{2}\\&y=-\dfrac{1}{2}}\Rightarrow H\left(-\dfrac{3}{2};-\dfrac{1}{2}\right).$$
		Vì $H$ là trung điểm của $BB'$ nên ta có\\ 
		$$\heva{&\dfrac{x_B+x_{B'}}{2}=x_H\\&\dfrac{y_B+y_{B'}}{2}=y_H}\Leftrightarrow  \heva{&x_{B'}=2x_H-x_B=2\cdot\left(-\dfrac{3}{2}\right)-0=-3\\&y_{B'}=2y_H-y_B=2\cdot\left(-\dfrac{1}{2}\right)-1=-2}\Rightarrow B'(-3;-2).$$
		Ta có  $\overrightarrow{AB'}=(6;3)\Rightarrow AB'=\sqrt{6^2+3^2}=3\sqrt{5}$. \\
		Vậy $a^2+2b=3^2+  10 =19$.
	} 
\end{bt}


%%%%%%%%%%%%%%%% TỰ LUẬN %%%%%%%%%%%%%%%%

% nguồn: MaHoaNH_TuLuan_Lop10
\begin{bt}%[0H9H3-6]%[ma:ba78aedd]
	Trong mặt phẳng tọa độ $Oxy$, cho  đường thẳng $d\colon \heva{&{x=-2-2t} \\
	&{y=1+2t}}$ và điểm $M(3;1)$. Xác định tung độ $M'$ đối xứng  với $M$ qua đường thẳng $d$.
	 \loigiai{
		Ta có  $d\colon\heva{&{x=-2-2t} \\&{y=1+2t}}\Rightarrow d\colon x+y+1=0$.\\
		Phương trình đường thẳng $MI$ là $x-y-2=0$.\\ 
		Toạ độ toạ hình chiếu $I$ của điểm $M$ là nghiệm hệ phương  trình
		\[\left\{\begin{array}{l}{x+y=-1} \\ {x-y=2}
	\end{array}\Leftrightarrow \left\{\begin{array}{l}{x=\dfrac{1}{2}} \\ {y=-\dfrac{3}{2}}
	\end{array}\Rightarrow I\left(\dfrac{1}{2};-\dfrac{3}{2} \right).\right. \right.\]
	Vì $M'$ đối xứng với $M$ qua đường thẳng $d$ nên $I$ là trung  điểm của $MM'$
	\[\Rightarrow \left\{\begin{array}{l}{x_{M'}=2x_I-x_M=-2} \\ {y_{M'}=2y_I-y_M=-4}
	\end{array}\Rightarrow M'(-2;-4)\right..\]
	Suy ra tung  độ $M{'}$ là $-4$.
	} 
\end{bt}

% nguồn: MaHoaNH_TuLuan_Lop10
\begin{bt}%[0H9H3-6]%[ma:b0d34620]
	Trong mặt phẳng  toạ độ $Oxy$, cho hình bình hành $ABCD$ có diện tích bằng $2$. Biết $A(0;2),B(3;0)$ và giao điểm $I$ của hai đường chéo hình bình hành nằm trên đường thẳng $y=-x$. Tìm hoành độ đỉnh $D$, biết $x_D>-14$.
	\loigiai{ 
		Phương trình đường thẳng $AB$ là $\dfrac{x}{3}+\dfrac{y}{2}=1\Leftrightarrow  2x+3y-6=0$.\\
		Điểm $I$ nằm trên đường thẳng $y=-x$ nên giả  sử $I(t;-t)$.\\
		 Vì $I$ là trung điểm của $AC$ nên $C(2t;-2t-2),I$ là trung điểm của $BD$ nên $D(2t-3;-2t)$.\\
		Ta có  $AB=\sqrt{(3-0)^2+(0-2)^2}=\sqrt{13}$.\\
		Suy ra  $\mathrm{d}(C,AB)=\dfrac{2}{\sqrt{13}}$.\\
		Khi đó $\dfrac{|2\cdot 2t+3(-2t-2)-6|}{\sqrt{2^2+3^2}}=\dfrac{2}{\sqrt{13}} \Leftrightarrow|-2t-12|=2$.\\ 
		Suy ra  $t=-5$ hoặc $t=-7$.
		\begin{itemize}
			\item Với $t=-5$, ta có  $C(-10;8),D(-13;10)$ nhận vì $x_D=-13>-14$.
			\item Với $t=-7$, ta có $C(-14;12),D(-17;14)$ loại  vì $x_D=-17<-14$.
		\end{itemize}
	} 
\end{bt}

% nguồn: MaHoaNH_TuLuan_Lop10
\begin{bt}%[0H9V3-6]%[ma:93e43c5e]
	Trong mặt phẳng $Oxy$, cho hai điểm $A(3;-3)$, $B(5;-1)$ và đường thẳng $\Delta: 2x - y - 1 = 0$. Tìm tổng hoành độ và tung độ của tất cả các điểm $M$ thuộc $\Delta$ sao  cho tam giác $MAB$ cân tại $M$.
	\loigiai{ 
		Đường thẳng $\Delta$ có phương trình tham  số
		$\heva{ & x = t \\ &  y = -1 + 2t.}$ \\
		Với  $M \in \Delta$ nên $M(t; -1 + 2t)$. \\
		Tam  giác $MAB$ cân tại $M$ nên $MA = MB$. \\
		$\Rightarrow (3 - t)^2 + (-2 - 2t)^2 = (5  - t)^2 + (-2t)^2 \Rightarrow 13 + 2t = 25 - 10t \Rightarrow 12t = 12 \Rightarrow t = 1$. \\
		Vậy điểm $M$ cân tại là $M(1;1) \Rightarrow  1 + 1 = 2$.
	} 
\end{bt}

% nguồn: MaHoaNH_TuLuan_Lop10
\begin{bt}%[0H9V3-6]%[ma:95f7fc1c]
	Trong mặt phẳng tọa độ $Oxy$, cho các điểm $A(-1; 3)$, $B(2;  6)$, $C(5; 0)$ và đường thẳng $\Delta\colon 3x-y+1=0$. Biết điểm $M(a; b)$ nằm trên $\Delta$ thì biểu thức $\left|\overrightarrow {MA}+\overrightarrow {MB}+\overrightarrow {MC}\right|+\left|\overrightarrow {MA}+2\overrightarrow {MB}\right|$ có giá trị nhỏ nhất. Tính giá trị của biểu thức $5a+10b$?
	\loigiai{ 
		Gọi $G$ là điểm thỏa mãn $\overrightarrow {GA}+\overrightarrow {GB}+\overrightarrow {GC}=0$, ta tính được tọa độ  điểm $G(2;3)$.\\
		Gọi $N$ là điểm  thỏa mãn $\overrightarrow {NA}+2\overrightarrow {NB}=0$, ta tính được tọa độ điểm $N(1;5)$.\\
		Từ đó ta thấy $G$, $N$ nằm về hai phía  so với đường thẳng $\Delta$.\\
		 Ta có $\left|\overrightarrow {MA}+\overrightarrow {MB}+\overrightarrow {MC}\right|=3\left|\overrightarrow {MG}\right|=3MG$ và $\left|\overrightarrow {MA}+2\overrightarrow {MB}\right|=3\left|\overrightarrow {MN}\right|=3MN$.\\
		Khi  đó $\left|\overrightarrow {MA}+\overrightarrow {MB}+\overrightarrow {MC}\right|+\left|\overrightarrow {MA}+2\overrightarrow {MB}\right|=3(MG+MN) \ge 3GN$.\\
		 Do đó $\left|\overrightarrow {MA}+\overrightarrow {MB}+\overrightarrow {MC}\right|+\left|\overrightarrow {MA}+2\overrightarrow {MB}\right|$ nhỏ nhất là bằng $3GN$, đạt được khi $3$ điểm $G$, $M$, $N$ thẳng hàng.\\
		 Suy ra $M$ là giao điểm của đường thẳng $GN$ và $\Delta$.\\
		Ta có $\overrightarrow {GN}=(-1;2)$, phương trình đường thẳng $GN$ là 
		$(x-1)+(y-5)=0 \Leftrightarrow 2x+y-7=0$.\\ 
		 Tọa độ điểm $M$ thoả $\heva{&2x+y-7=0 \\ &3x-y+1=0}\Leftrightarrow \heva {&x=\dfrac{6}{5} \\ &y=\dfrac{23}{5}.}$\\
		Vậy $5a+10b=35$. 
	} 
\end{bt}
